\section{Normal local cohomology and the first fusion relation}
\label{sec:normal-locality}

Leading logarithmic fusion remembers normal directions through a
twisted differential.  To retain the local cohomology class of a pole,
we instead use the interval algebra with spatial differential zero.
Completion preserves its finite principal parts.  The resulting
three-input comparison detects an obstruction that disappears when
the normal parameter is inverted.

\subsection{A relative interval containing the first pole}

Let \(F\) be a commutative \(\mathbb Q\)-algebra. Put
\(R=F[t]\), \(B=F[[t]]\), and \(L=B[t^{-1}]\).
Series in \(L\) have a finite lower bound on their exponents.
All spatial variables have degree zero and differential zero.
Complexes are cohomological. Tensor products use the Koszul sign.

For a ring \(T\) in which \(t\) is a nonzerodivisor, write
\begin{equation}\label{nl:eq:interval}
 \A(T)=\mathcal I_T(t)
\end{equation}
for the interval algebra of \eqref{sb:eq:general-interval}.
Its differential is denoted by \(\dco\); its degree-zero module
is \(Te_0\oplus Te_1\), and its overlap module is
\(T[t^{-1}]\theta\).  This notation retains the first-endpoint
augmentation and the two distinct endpoint actions.

The element \(\theta/t\) is closed. A primitive would require
an endpoint coefficient with a negative power of \(t\).
Such a coefficient is absent from \(\A(B)^0\).

\begin{theorem}[Principal parts and completion]\label{nl:thm:relative}
The augmentation ideal of \(\A(T)\) is the complex
\[
 \overline\A(T)=(Te_1\xrightarrow{\dco}T[t^{-1}]\theta),
 \qquad \dco(fe_1)=f\theta,
\]
in degrees zero and one. Thus
\[
 H^0(\A(T))=T,\qquad
 H^1(\A(T))=T[t^{-1}]/T.
\]
The map \(\overline\A(R)\to\overline\A(B)\) induced by
power-series expansion is a quasi-isomorphism of \(R\)-complexes.
On degree-one cohomology it is the isomorphism
\begin{equation}\label{nl:eq:principal-parts}
 F[t,t^{-1}]/F[t]\xrightarrow{\ \sim\ }L/B
      =\bigoplus_{j\geq1}Ft^{-j}.
\end{equation}
The full interval map has the different degree-zero map \(R\to B\).
If \(T\) is Noetherian, \(\overline\A(T)\) represents
\(R\Gamma_{(t)}(T)\).
\end{theorem}

\begin{proof}
The augmentation splits by the unit, as a map of complexes.
On its kernel the differential is the inclusion into the localization.
This proves the cohomology assertions. Every element of \(L\)
has a unique finite negative part modulo \(B\). This proves
\eqref{nl:eq:principal-parts}, including injectivity and surjectivity.
The local cohomology identification is the one-generator extended
\v{C}ech complex \cite{stacks-local}*{Lemma 47.10.2}.
\end{proof}

For a fixed integer \(N\geq0\), the subcomplex with degree-one
module \(t^{-N}B\theta\) has cohomology \(t^{-N}B/B\).
For \(N>0\), this module contains \([\theta/t]\).
No spatial de Rham differential or logarithmic differential is applied.

\subsection{Normal substitution and actual coefficient bars}

Let \(T_0\) be a commutative ring, and let
\(\delta_1,\ldots,\delta_m\) be elements of \(T_0\).
The interval \(\A_{\delta_i}(T_0)\) is defined by
\eqref{nl:eq:interval}, with \(\delta_i\) in place of \(t\).
Let \(\phi:T_0\to B\) satisfy
\[
 \phi(\delta_i)=t^{n_i}v_i,\qquad n_i\geq0,\quad v_i\in B^\times.
\]
Localization and scalar extension give a strict augmented algebra map
\begin{equation}\label{nl:eq:substitution}
 \bigotimes_{i=1}^m\A_{\delta_i}(T_0)
   \longrightarrow
 \bigotimes_{i=1}^m\A_{t^{n_i}v_i}(B).
\end{equation}
It applies \(\phi\) to each coefficient and preserves each endpoint
and overlap symbol. Every tensor product is over its displayed base ring.
If \(n_i>0\), the \(i\)-th target factor is \(\A(B)\).
If \(n_i=0\), its degree-one module is \(B\), and its augmentation
ideal contracts by \(h(f\theta)=fe_1\).
There is a canonical isomorphism between the target and scalar extension
of the source. This statement requires no flatness of \(\phi\).
It makes no assertion that \eqref{nl:eq:substitution} is a quasi-isomorphism.

Let \(D\) be an augmented differential graded algebra over \(B\).
Its reduced bar is
\[
 \mathrm B_B(D)=\bigoplus_{r\geq0}(s\overline D)^{\otimes_B r},
 \qquad |sa|=|a|-1.
\]
On one and two entries its degree-one differential has components
\begin{equation}\label{nl:eq:bar}
 b_1[a]=[\dco a],\qquad
 b_2[a|b]=(-1)^{|a|+1}[ab].
\end{equation}
An operation after the first \(j-1\) entries has the prefix sign
\((-1)^{\sum_{i<j}(|a_i|-1)}\). Cuts split a word into consecutive
subwords, including either empty subword.

\begin{proposition}[Coefficient bar substitution]\label{nl:prop:bar}
Formula~\eqref{nl:eq:bar} squares to zero and commutes with cuts.
After scalar extension, \eqref{nl:eq:substitution} induces a strict
coalgebra chain map on these bars. The same assertion holds after
tensoring with an augmented state algebra and applying its identity map.
All words and input state polynomials are finite.
\end{proposition}

\begin{proof}
On one entry, the square is \(\dco^2\). On two entries,
the mixed terms cancel by the Leibniz identity. On three entries,
the two multiplication terms are
\((-1)^{|b|}[(ab)c]\) and \(-(-1)^{|b|}[a(bc)]\).
They cancel by associativity. Disjoint operations cancel because
each has degree one. A cut commutes with every operation confined
to one side. A multiplication across a proposed cut removes that
cut on both sides of the chain equation. This proves the first assertion.
Apply the algebra map to each entry for the second. It preserves
the augmentation, the internal differential, multiplication, and cuts.
Scalar extension commutes with the finite tensor powers and their direct sum.
\end{proof}

For example, \([\theta/t]\) is a closed length-one element of the
normal coefficient bar. It remains a nonzero class in its length-one
subcomplex. This assertion does not identify its class in the full bar.
The multiplication in Proposition~\ref{nl:prop:bar} is the displayed
coefficient and state algebra multiplication. The collision map defined
next is a separate operation with its own ternary equation.

\subsection{Binary contraction and the three-point normal chart}

Use the one-species state algebra \(\V\) of
\eqref{fs:eq:state-algebra} over \(\mathbb Q\), with its auxiliary
weight and Hasse translation.  In this section
\(D_{ij}=\mathsf D_{ij}\) is the rational odd contraction
\eqref{fs:eq:rational-pair-operator}.  Coefficients precede states,
so moving \(D_{ij}\) past an overlap symbol changes its sign.

At separation \(v\), with \(v=t\) times a unit of \(B\), define
\begin{equation}\label{nl:eq:binary}
 Y_v=E_v(1+\theta_vD_{12}),\qquad E_v(a\otimes b)=(\tau_va)b.
\end{equation}
The uncontracted term has equal endpoint values. The contraction term
lies in degree one of the coefficient complex. Both are closed,
and the contraction has total degree zero. Thus \(Y_v\) is a
degree-zero chain map on polynomial inputs.
For each input it takes values in
\[
 \bigcup_{N\geq0}\prod_{w\geq0}
 F_N(\A(B)\otimes_{\mathbb Q}\V_w),
\]
where degree-zero coefficient modules are \(B^2\), and positive
coefficient degrees have bound \(t^{-N}B\).
Only finitely many derivatives occur before translation. Translation
adds nonnegative powers of \(t\), so one bound works for every output weight.
In particular,
\[
 Y_t(x_0,\eta_0)=((\tau_tx_0)\eta_0)\,1+\frac{\theta_t}{t}\,1_{\V}.
\]
The scalar-state part represents the first nonzero local cohomology class.

For the total three-point collision, use the coordinate ring and chart
\begin{equation}\label{nl:eq:normal-chart}
 F=\mathbb Q[u,u^{-1},(1-u)^{-1}],\quad B=F[[t]],\qquad
 z_1=t,\quad z_2=t(1-u),\quad z_3=0.
\end{equation}
Thus
\(a=z_1-z_2=tu\), \(b=z_2-z_3=t(1-u)\), and
\(c=z_1-z_3=t\). All three pair intervals have overlap ring \(L\).
Retain their distinct symbols and put
\[
 E_L=\A_a(B)\otimes_B\A_b(B),\qquad
 E_R=\A_b(B)\otimes_B\A_c(B),\qquad
 E(\alpha\otimes\beta\otimes\gamma)
       =(\tau_c\alpha)(\tau_b\beta)\gamma.
\]

\begin{proposition}[The actual two iterations]\label{nl:prop:iterations}
The rational representatives of the binary iterations are
\begin{equation}\label{nl:eq:iterations}
 \begin{aligned}
 Y_L&=E(1+\theta_b(D_{13}+D_{23}))(1+\theta_aD_{12}),\\
 Y_R&=E(1+\theta_c(D_{12}+D_{13}))(1+\theta_bD_{23}).
 \end{aligned}
\end{equation}
They extend to degree-zero chain maps
\(B\otimes_{\mathbb Q}\V^{\otimes3}\to
\bigcup_N\prod_w F_N(E_L\otimes_{\mathbb Q}\V_w)\)
and to the analogous target for \(E_R\).
The source has differential zero. Each formula has at most two contractions.
\end{proposition}

\begin{proof}
For the left iteration, the inner contraction is \(D_{12}\).
At the outer merger, differentiation of a product sums the contributions
from its two original factors. It gives \(D_{13}+D_{23}\).
The right iteration similarly gives \(D_{23}\) and \(D_{12}+D_{13}\).
The original leaf separations follow from the formal identities
\begin{align*}
 \sum_{j\geq0}\binom{r+j}{r}v^jK_{r+j,s}(w)&=K_{rs}(w+v),\\
 \sum_{j\geq0}\binom{s+j}{s}v^jK_{r,s+j}(w)&=K_{rs}(w-v).
\end{align*}
Expand the right sides in nonnegative powers of \(v/w\).
The negative-binomial formula proves the identities term by term.
Together with the translation law, they give \eqref{nl:eq:iterations}
in each iterated expansion domain before normal substitution.

Under \eqref{nl:eq:normal-chart}, every denominator is a power of
\(t\) times a unit of \(F\). Every contraction carries an
overlap symbol, whose module permits that power of \(t\).
Products of two overlaps permit the product denominator.
The endpoint terms are regular. Each coefficient module embeds in
its localization at \(t\), so this extension of the rational representative
is unique. Finiteness before translation proves the uniform weight bound.
All terms are closed with the displayed total degree.
\end{proof}

The iterated Laurent targets used to derive these formulas are distinct.
Equality of their rational representatives supplies the normal expressions.
No operation on arbitrary completed intermediate states is asserted.
In particular, extension of these particular state maps across \(t=0\)
does not extend the endpoint ring map \(\mathbb Q[a,b,c^{-1}]\to B\).
That map would require \(t^{-1}\in B\).

\subsection{The ternary obstruction and its exact localization}

There is a strict augmented algebra isomorphism
\[
 r:E_L\longrightarrow E_R
\]
which relabels the \(a\)-interval as the \(c\)-interval, preserves
the \(b\)-interval, and uses the signed interchange of the two factors.
This is possible because their overlap rings are all \(L\).
It fixes every element of \(B\).
For a degree-minus-one map \(H\) from the state source to the
right output, the first ternary equation would be
\begin{equation}\label{nl:eq:ternary}
 d_{\mathrm{Hom}}H=\dco H+H\dco_{\V^{\otimes3}}
                  =(r\otimes1)Y_L-Y_R.
\end{equation}

\begin{theorem}[First pole obstruction]\label{nl:thm:obstruction}
Equation~\eqref{nl:eq:ternary} has no solution over \(B\).
More generally, there is no \(B\)-linear coefficient chain map
\(q:E_L\to E_R\) for which \((q\otimes1)Y_L\) and \(Y_R\)
are chain homotopic on all polynomial state inputs.
This includes coefficient maps that need not be algebra maps or equivalences.
The state factor in this assertion is the identity.
\end{theorem}

\begin{proof}
For inputs \((x_0,1,\eta_0)\), take the scalar-state coefficient.
Only \(D_{13}\) acts, with coefficient \(1/c\). Hence
\begin{equation}\label{nl:eq:deciding-pole}
 [((r\otimes1)Y_L-Y_R)(x_0,1,\eta_0)]_{1_{\V}}
       =\frac{\theta_b}{t}-\frac{\theta_c}{t}.
\end{equation}
First-endpoint evaluation on the \(c\)-factor is a chain map
\(E_R\to\A_b(B)\). It takes this difference to \(\theta_b/t\),
which is not a boundary by Theorem~\ref{nl:thm:relative}.
Scalar-state projection is also a chain map, since the state differential
is zero. Thus no primitive of the full output exists.

For the stronger assertion, use also \((1,x_0,\eta_0)\).
Its two scalar outputs are \(\theta_b/b\).
An intertwining homotopy would therefore imply
\[
 q_*[\theta_b/b]=[\theta_b/b],\qquad
 q_*[\theta_b/c]=[\theta_c/c]
\]
in the corresponding degree-one coefficient cohomology groups.
But \(b/c=1-u\in B^\times\), so \(B\)-linearity forces
\(q_*[\theta_b/c]=[\theta_b/c]\).
The two required values differ by the nonzero class
\eqref{nl:eq:deciding-pole}. This is a contradiction.
\end{proof}

\begin{theorem}[The exact localization for the ternary relation]
\label{nl:thm:localization}
Among nonzero localizations \(B'\) of \(B\), equation
\eqref{nl:eq:ternary} has a solution for all polynomial inputs if and
only if \(t\) is invertible in \(B'\).
Consequently \(B[t^{-1}]\) is the least such endpoint localization.
It removes the entire principal local cohomology module \(L/B\).
\end{theorem}

\begin{proof}
If \(t\) is not invertible, \(1/t\notin B'\), so
\([1/t]\ne0\) in \(B'[t^{-1}]/B'\).
The proof of \eqref{nl:eq:deciding-pole} applies over \(B'\).
This proves necessity without a flatness or generic-point argument.

Suppose \(t\) is invertible. For each interval define
\[
 h_i(f\theta_i)=fe_1^i,\qquad h_i(\A_i^0)=0,\qquad
 F_i=1_i\varepsilon_i.
\]
These maps are legal because every overlap coefficient is now an endpoint
coefficient. They satisfy \(\dco h_i+h_i\dco=1-F_i\).
On \(E_R\), set
\[
 h_R=h_b\otimes1+F_b\otimes h_c,
 \qquad Z=(r\otimes1)Y_L-Y_R,
 \qquad H=(h_R\otimes1)Z.
\]
The tensor notation includes its Koszul signs. Differentiation gives
\(\dco h_R+h_R\dco=1-F_b\otimes F_c\).
The map \(Z\) is closed. Full endpoint evaluation kills every contraction,
and the two uncontracted terms both equal \(E\). Thus
\((F_b\otimes F_c\otimes1)Z=0\), and \(\dco H=Z\).
This is \eqref{nl:eq:ternary}, including its source differential.
Finally \(B'[t^{-1}]=B'\), so its principal local cohomology vanishes.
\end{proof}

For the deciding input, one obtains the explicit scalar primitive
\[
 [H(x_0,1,\eta_0)]_{1_{\V}}
       =\frac{e_1^b}{t}-\frac{e_1^c}{t}.
\]
It is legal exactly on the stated localization. A ternary correction
on a suspended collision word would require the same unsuspended Hom
equation, with the chosen suspension sign. Higher operations cannot
supply a missing primitive in that equation.

\subsection{A shared interval and the cohomology it removes}

There is a smaller common target over \(B\) itself. Let
\(q_L:E_L\to\A(B)\) relabel both factors as \(\A(B)\) and
multiply them in the order \(a,b\). Let \(q_R:E_R\to\A(B)\)
use the order \(b,c\). Multiplication is a degree-zero chain map
from the tensor product of complexes, so these maps are well defined.
They preserve the unit and augmentation.

\begin{proposition}[Equality after interval multiplication]\label{nl:prop:quotient}
On every polynomial input,
\begin{equation}\label{nl:eq:shared}
 (q_L\otimes1)Y_L=(q_R\otimes1)Y_R
       =E(1+\theta(D_{12}+D_{13}+D_{23})).
\end{equation}
The target retains \([\theta/t]\ne0\).
Neither \(q_L\) nor \(q_R\) is a quasi-isomorphism or an algebra map.
Applying either map to each entry fails to give a strict chain map
on the reduced coefficient bars.
\end{proposition}

\begin{proof}
Both overlap symbols map to \(\theta\), while their product maps to
\(\theta^2=0\). The degree-zero coefficients in the state formulas
are diagonal scalars. Applying the maps to \eqref{nl:eq:iterations} gives
\eqref{nl:eq:shared} term by term, with all double contractions killed.
The pole assertion is Theorem~\ref{nl:thm:relative}.

For nonmultiplicativity, write \(q\) for multiplication of two copies
of \(\A(B)\). Take \(\alpha=1\otimes\theta\) and
\(\beta=e_1\otimes1\). Then
\[
 q(\alpha\beta)=q(e_1\otimes\theta)=0,\qquad
 q(\alpha)q(\beta)=\theta e_1=\theta.
\]
Both elements lie in the augmentation ideal. Since \(|\alpha|=1\),
the length-one component of the bar chain-map defect on
\([\alpha|\beta]\) is \([\theta]\), by \eqref{nl:eq:bar}.
The same example with \(\theta/t\) has defect \([\theta/t]\).
Internal differentials match because \(q\) is a chain map.
The failure of quasi-isomorphism follows from the next calculation.
\end{proof}

\begin{lemma}[Independent interval classes]\label{nl:lem:cube}
For \(m\geq1\), put \(Q_m=\A(B)^{\otimes_B m}\) and \(M=L/B\).
Then
\[
 H^0(Q_m)=B,\qquad H^1(Q_m)\cong M^{\oplus(2^m-1)},\qquad
 H^j(Q_m)=0\quad(j\geq2).
\]
For \(m=2\), the multiplication map on cohomology is a split
surjection \(M^3\to M\) with kernel \(M^2\).
\end{lemma}

\begin{proof}
In positive coefficient degree, every summand of \(Q_m\) is a copy
of \(L\), since \(L\otimes_B L=L\). Its differential is the
oriented cubical cochain differential. In degree zero, the module is
\(B^{2^m}\). The fully localized cubical complex has degree-zero
kernel the diagonal \(L\) and vanishing higher cohomology.
An explicit contraction follows by tensoring the single-interval
endpoint contractions used in Theorem~\ref{nl:thm:localization}.
The differentials in positive degrees are unchanged before localization.
Consequently
\[
 H^1(Q_m)=L^{2^m}/(B^{2^m}+\operatorname{diag}L)
              \cong M^{2^m}/\operatorname{diag}M.
\]
Subtract the coordinate at one fixed vertex to obtain \(M^{2^m-1}\).
The degree-zero kernel is \(\operatorname{diag}B\), and all higher
groups vanish by the localized contraction.

For \(m=2\), multiplication selects the vertices \(00\) and \(11\)
in degree zero. In the preceding presentation it takes a vertex vector
to the difference of its \(11\) and \(00\) coordinates modulo \(B\).
After subtracting the \(00\) coordinate, this is projection \(M^3\to M\).
It has kernel \(M^2\). Inserting a unit into one interval gives a
chain section and the stated splitting.
\end{proof}

The normal coefficient bars therefore exist with their original
endpoint differential and finite words. The binary contraction retains
its pole class. The native three-input relation fails in independent
interval carriers even after their coefficient rings coincide.
The shared-interval map gives the exact equality \eqref{nl:eq:shared},
with the explicit loss \(M^2\) and the displayed strict bar defect.
An action through that target would require higher bar comparison maps
and a domain for repeated operations on completed state outputs.
Those constructions are additional to the coefficient bars and the
binary and ternary formulas proved here.

The underlying principal-part module nevertheless has a normal
derivative description.  Section~\ref{ns:sec} constructs this
description and its supported Koszul comparison, keeping the
module maps separate from the coefficient-algebra and bar maps.
